\documentclass[11pt]{article}
\usepackage[margin=1in]{geometry}
\usepackage{parskip}
\usepackage{enumitem}
\usepackage[hidelinks]{hyperref}
\setlist[itemize]{leftmargin=1.4em, topsep=2pt, itemsep=1pt}
\pagestyle{empty}

\begin{document}

\noindent\today

\noindent To the Editor-in-Chief and the Editorial Board,\\
\textit{IEEE Transactions on Network and Service Management}

\bigskip
\noindent\textbf{Subject:} Submission of an original manuscript for consideration as a
regular paper.

\bigskip
\noindent Dear Editor,

We are pleased to submit our manuscript, ``\textbf{Learning to Route Under Load: An
Inductive Graph Network that Amortizes Congestion-Aware Traffic Engineering for LEO
Mega-Constellations},'' for consideration for publication in the \emph{IEEE Transactions
on Network and Service Management}.

\textbf{Problem and contribution.} Low Earth orbit (LEO) mega-constellations route traffic
over a large, fast-changing mesh of inter-satellite links whose topology is known in
advance from ephemeris. Under load, routing on propagation delay alone overloads a few
bottleneck links while alternatives sit idle, inflating end-to-end travel time by up to a
factor of forty in our study. The congestion-aware optimum, the user equilibrium (UE),
removes this waste but requires the full demand matrix and an iterative solve every slot,
which does not scale to constellations of more than a thousand satellites with
seconds-scale topology change. We formalize routing under load as the \emph{amortization}
of a traffic-assignment equilibrium and solve it with an inductive graph neural network
that predicts the equilibrium congestion price on every link in a single forward pass from
cheap, observable load features, decoded by a one-shot multipath router. The method:
\begin{itemize}
  \item closely approaches UE in distribution (recovering about $97\%$ of the gain over
        blind routing);
  \item transfers \emph{zero shot} to a shell twelve times larger, still recovering
        $94\pm2\%$ of the gain, as the best deployable policy;
  \item runs $21$ to $29\times$ faster than the iterative solve, with the saving growing
        with constellation size;
  \item supports proactive routing under drifting traffic hotspots, where a reactive
        controller degrades below congestion-blind routing.
\end{itemize}
We additionally contribute a proposition bounding the travel-time suboptimality of the
decoded routing by the price-prediction error (validated empirically), and a documented set
of evaluation pitfalls specific to learned traffic engineering, with a protocol to avoid
them. We also report, transparently, a negative result: a quantum-walk propagation operator
that motivated the study offers no benefit once load features are available.

\textbf{Fit to the journal.} The work lies squarely within the scope of \emph{IEEE TNSM}:
it concerns network and service management for non-terrestrial networks, specifically
congestion-aware traffic engineering, learning-based network control, and resource
management at scale. It pairs a management-oriented problem formulation with a deployable,
inexpensive control mechanism and an honest, reproducible evaluation methodology.

\textbf{Reproducibility.} The complete simulator, models, experiments, and
figure-generation scripts are publicly released; every reported number is regenerated from
result files by an included script. All results follow from constellation parameters with
no external data dependency.

\textbf{Declarations.} This manuscript is original, has not been published previously, and
is not under consideration for publication elsewhere, in whole or in part. All authors have
read and approved the submission and agree to its content. The authors declare no conflict
of interest. This research is funded by the Science and Technology Development Fund of The
University of Danang under project number B2025-DN02-26.

Thank you for considering our submission. We look forward to the reviewers' feedback.

\bigskip
\noindent Sincerely,\\
Phuc Hao Do, Nang Hung Van Nguyen, and Minh Tuan Pham

\bigskip
\noindent\textbf{Corresponding author:}\\
Minh Tuan Pham\\
Faculty of Information Technology,\\
The University of Danang -- University of Science and Technology, Da Nang, Viet Nam\\
E-mail: \href{mailto:pmtuan@dut.udn.vn}{pmtuan@dut.udn.vn}

\end{document}
